\documentclass[11pt,a4paper]{article}
\usepackage[margin=25mm,headheight=14pt]{geometry}
\usepackage[T1]{fontenc}
\usepackage[utf8]{inputenc}
\usepackage{lmodern,amsmath,amssymb}
\usepackage{xcolor,graphicx,booktabs,longtable,array}
\usepackage{microtype,enumitem,fancyhdr,titlesec}
\usepackage{hyperref}
\definecolor{accent}{HTML}{163B55}
\hypersetup{colorlinks=true,linkcolor=accent,urlcolor=accent,
 pdftitle={Dense nonperiodic limit-periodic potentials with finite interval spectra},pdfauthor={Matthew J. Colbrook},
 pdfsubject={AIM 1031: mathematical result and proof}}
\urlstyle{same}
\Urlmuskip=0mu plus 1mu
\setlength{\parindent}{0pt}
\setlength{\parskip}{0.35em}
\setlength{\emergencystretch}{3em}
\setlist{itemsep=0.25em,topsep=0.3em}
\providecommand{\tightlist}{\setlength{\itemsep}{0pt}\setlength{\parskip}{0pt}}
\titleformat{\section}{\normalfont\large\bfseries\color{accent}}{\thesection}{0.7em}{}
\titleformat{\subsection}{\normalfont\normalsize\bfseries}{\thesubsection}{0.6em}{}
\pagestyle{fancy}
\fancyhf{}
\fancyhead[L]{\small\sffamily AIM 1031}
\fancyhead[R]{\small\sffamily Matthew J. Colbrook}
\fancyfoot[C]{\thepage}
\renewcommand{\headrulewidth}{0.3pt}
\setcounter{secnumdepth}{0}
\allowdisplaybreaks[2]
\title{\textbf{Dense nonperiodic limit-periodic potentials with finite interval spectra}}
\author{Matthew J. Colbrook\\[3pt]
{\small Department of Applied Mathematics and Theoretical Physics}\\
{\small University of Cambridge, United Kingdom}\\
{\small\href{mailto:mjc249@maths.cam.ac.uk}{\texttt{mjc249@maths.cam.ac.uk}}}}
\date{3 November 2026}
\begin{document}
\maketitle
\begin{abstract}
We prove that real nonperiodic limit-periodic potentials whose spectrum is a finite union of nondegenerate intervals are dense among all real limit-periodic potentials. Finite block insertions fill limiting spectral clusters, while separately protected marker blocks ensure nonperiodicity. This answers AIM 1031 affirmatively.
\end{abstract}
\noindent\textbf{Repository target.}
This manuscript concerns AIM problem 1031 at revision
\texttt{3fa50884d129}. Its scope is stated below.
\medskip

The answer to Problem 1031 is \textbf{yes}. In fact, potentials whose
entire spectrum is a finite union of nondegenerate closed intervals are
dense among limit-periodic potentials, and the approximating potentials
can always be required to be nonperiodic. The proof places
finite-support approximate eigenvectors inside long periodic blocks.
Small constant shifts of those blocks fill successively smaller spectral
gaps, while separate marker blocks ensure that the limiting potential is
nonperiodic.

\hypertarget{1031-theorem}{%
\section{Theorem}\label{1031-theorem}}

Let \(V:\mathbb Z\to\mathbb R\) be a bounded uniform limit of periodic
sequences, and let \(\varepsilon>0\). Then there exists a nonperiodic
uniform limit of periodic sequences \(W\) such that

\[
\|W-V\|_\infty<\varepsilon,
\qquad
\sigma(H_W)=\bigcup_{j=1}^{M}[\ell_j,r_j],
\qquad \ell_j<r_j<\ell_{j+1},
\]

where

\[
(H_Wu)(n)=u(n-1)+u(n+1)+W(n)u(n)
\quad\text{on }\ell^2(\mathbb Z).
\]

The interval union is finite, with \(M\geq1\), and the displayed
separation condition applies for \(j<M\). In particular, the spectrum
has nonempty interior and is not a Cantor set. The initial potential
need not itself be nonperiodic. Thus the theorem has all the quantifiers
required in the catalogue question, with a stronger conclusion. If \(V\)
is constant, one may take \(M=1\).

\hypertarget{1031-a-finite-block-insertion-lemma}{%
\section{1. A finite block insertion
lemma}\label{1031-a-finite-block-insertion-lemma}}

Let \(P\) be a bounded real potential of integer period \(q\). Fix a
number \(a>0\), a positive integer \(K\), and a finite list of target
energies \(E_1,\ldots,E_m\). Suppose that

\[
\operatorname{dist}(E_j,\sigma(H_P))\leq a
\qquad(1\leq j\leq m).
\]

For every \(\eta>0\), there is a periodic potential \(\widetilde P\),
whose period is a multiple of \(q\), with the following three
properties:

\[
\|\widetilde P-P\|_\infty\leq a,
\tag{1}
\]

\[
\operatorname{dist}(E_j,\sigma(H_{\widetilde P}))<\eta
\qquad(1\leq j\leq m),
\tag{2}
\]

and, for every integer \(1\leq p\leq K\), there is an integer \(n_p\)
such that

\[
|\widetilde P(n_p+p)-\widetilde P(n_p)|\geq2a.
\tag{3}
\]

\hypertarget{1031-proof-of-the-lemma}{%
\subsection{Proof of the lemma}\label{1031-proof-of-the-lemma}}

Choose \(e_j\in\sigma(H_P)\) with \(|E_j-e_j|\leq a\), and set
\(c_j=E_j-e_j\). Since adding a constant to a potential shifts its
spectrum by that constant,

\[
E_j\in\sigma(H_{P+c_j}).
\]

For a bounded self-adjoint operator, every spectral point is an
approximate eigenvalue. Finitely supported sequences are dense in
\(\ell^2(\mathbb Z)\), and the operator is bounded, so there is a
finitely supported unit vector \(\phi_j\) with

\[
\|(H_{P+c_j}-E_j)\phi_j\|_2<\eta.
\tag{4}
\]

Choose an integer block length \(L\), divisible by \(q\), which exceeds
\(K\) and is sufficiently large to accommodate each support after
translation by a multiple of \(q\). On a finite sequence of consecutive
blocks of length \(L\), prescribe the potential as follows:

\[
P+c_1,\ P+c_2,\ \ldots,\ P+c_m,\ P+a,\ P-a,\ P+a.
\tag{5}
\]

Then repeat this entire block pattern periodically. Because every block
boundary is a multiple of \(q\), the background \(P\) has the same phase
on every block. The resulting potential \(\widetilde P\) has period
\((m+3)L\), a multiple of \(q\), and satisfies (1).

Translate \(\phi_j\) by a multiple of \(q\) so that its support lies in
the block carrying \(P+c_j\). On that support the two potentials agree
exactly. Their difference is a multiplication operator, so the
translated vector still has the residual in (4), even though the
potentials outside the block may differ. The spectral theorem gives

\[
\operatorname{dist}(E_j,\sigma(H_{\widetilde P}))
\leq\|(H_{\widetilde P}-E_j)\phi_j^{\rm translated}\|_2<\eta.
\]

This proves (2).

For (3), use the two internal boundaries in the three marker blocks
\(P+a,P-a,P+a\). Call them \(b_1,b_2\); both are multiples of \(q\). For
\(1\leq p\leq K<L\), the pair of sites \(b_i-1,b_i+p-1\) crosses the
corresponding boundary and stays within its two adjacent marker blocks.
In both cases the background difference is the same number

\[
v_p=P(p-1)-P(-1).
\]

The two differences of the new potential are therefore

\[
v_p-2a,\qquad v_p+2a.
\]

At least one has absolute value at least \(2a\). This proves (3)
simultaneously for every \(p\leq K\). \(\square\)

\hypertarget{1031-spectral-clusters-and-the-perturbation-budget}{%
\section{2. Spectral clusters and the perturbation
budget}\label{1031-spectral-clusters-and-the-perturbation-budget}}

Choose a periodic potential \(V_0\) with

\[
\|V-V_0\|_\infty<\frac{\varepsilon}{2}.
\]

Write its spectrum as its finitely many connected components:

\[
S_0=\sigma(H_{V_0})
=\bigcup_{j=1}^{M}I_{0,j},
\qquad
I_{0,j}=[\ell_{0,j},r_{0,j}],
\qquad
\ell_{0,j}<r_{0,j}<\ell_{0,j+1}.
\tag{6}
\]

The intervals are nondegenerate by the elementary Floquet description of
the spectrum of a periodic discrete Schrödinger operator. Touching bands
are combined into a single component. If \(M>1\), let

\[
G=\min_{j<M}(\ell_{0,j+1}-r_{0,j})>0.
\]

Choose a number \(t\) satisfying

\[
0<t<\varepsilon,
\qquad
 t<G/4\quad\hbox{if }M>1.
\]

For \(k\geq1\), set

\[
a_k=t4^{-k},
\qquad
\delta_k=\frac{a_{k+1}}8=\frac{a_k}{32},
\qquad
\delta_0=0.
\tag{7}
\]

Thus

\[
\sum_{k\geq1}a_k=\frac t3,
\qquad
\sum_{j>k}a_j=\frac{a_k}{3},
\qquad
\delta_{k-1}=\frac{a_k}{8}\quad(k\geq2).
\tag{8}
\]

The initial gaps will remain open throughout the construction. Indeed,
each endpoint can move by at most \(\sum a_k=t/3\), so all gaps between
the original \(M\) clusters remain at least \(G-2t/3>G/2\) when \(M>1\).

We shall construct periodic potentials \(V_k\), with periods divisible
by those of their predecessors, and ordered disjoint intervals
\(I_{k,1},\ldots,I_{k,M}\) such that, with \(S_k=\sigma(H_{V_k})\),

\[
\|V_k-V_{k-1}\|_\infty\leq a_k,
\tag{9}
\]

\[
S_k\subset\bigcup_{j=1}^{M}I_{k,j},
\qquad
\sup_{E\in\bigcup_jI_{k,j}}\operatorname{dist}(E,S_k)\leq\delta_k,
\tag{10}
\]

and each \(I_{k,j}\) is the convex hull of the part of \(S_k\) in its
corresponding cluster. Its endpoints therefore belong to \(S_k\). We
will also retain the marker condition

\[
\forall\,1\leq p\leq k\quad
\exists\,n_{k,p}\in\mathbb Z:\quad
|V_k(n_{k,p}+p)-V_k(n_{k,p})|\geq2a_k.
\tag{11}
\]

At \(k=0\), (10) holds with \(\delta_0=0\) by (6).

\hypertarget{1031-filling-each-expanded-cluster}{%
\section{3. Filling each expanded
cluster}\label{1031-filling-each-expanded-cluster}}

Suppose \(V_{k-1}\) and its cluster hulls have been constructed. Define
the expanded intervals

\[
K_{k,j}=I_{k-1,j}+[-a_k,a_k],
\qquad
K_k=\bigcup_{j=1}^{M}K_{k,j}.
\tag{12}
\]

The mesh bound in (10), together with \(\delta_{k-1}<a_k\), gives the
exact identity

\[
S_{k-1}+[-a_k,a_k]=K_k.
\tag{13}
\]

For completeness, a point inside a previous hull is within
\(\delta_{k-1}\) of \(S_{k-1}\); a point in one of its two extensions is
within \(a_k\) of the corresponding endpoint, which belongs to
\(S_{k-1}\). This proves one inclusion. The reverse inclusion follows
from \(S_{k-1}\subset\bigcup_j I_{k-1,j}\). The case \(k=1\) uses
\(\delta_0=0\).

Choose a finite \(\delta_k/4\)-net \(E_1,\ldots,E_m\) of the entire
compact set \(K_k\). By (13), every target satisfies

\[
\operatorname{dist}(E_j,S_{k-1})\leq a_k.
\]

Apply the finite block insertion lemma with \(P=V_{k-1}\), \(a=a_k\),
\(K=k\), and \(\eta=\delta_k/4\). It gives a periodic potential \(V_k\)
satisfying (9), (11), and

\[
\sup_{E\in K_k}\operatorname{dist}(E,S_k)\leq\delta_k.
\tag{14}
\]

The residual estimate actually gives \(\delta_k/2\), but the weaker
bound is convenient.

For bounded self-adjoint operators, spectral Hausdorff distance is at
most operator-norm distance. Since the difference of the Schrödinger
operators is multiplication by the difference of their potentials, (9)
and (13) also give the exact confinement

\[
S_k\subset S_{k-1}+[-a_k,a_k]=K_k.
\tag{15}
\]

This is what controls the entire spectrum, including spectral points
associated with interfaces between the blocks.

Define

\[
I_{k,j}=\operatorname{conv}(S_k\cap K_{k,j})
=[\ell_{k,j},r_{k,j}].
\tag{16}
\]

The sets being convexified are nonempty. If there is only one cluster,
this follows from (14). If \(M>1\), the gaps between the \(K_{k,j}\)
exceed \(G/2\), while \(\delta_k\leq t/128<G/512\); hence the spectral
points within \(\delta_k\) of either endpoint of \(K_{k,j}\) must lie in
that same cluster. Compactness then gives the endpoint bounds

\[
\ell_{k-1,j}-a_k
\leq\ell_{k,j}
\leq\ell_{k-1,j}-a_k+\delta_k,
\tag{17}
\]

\[
r_{k-1,j}+a_k-\delta_k
\leq r_{k,j}
\leq r_{k-1,j}+a_k.
\tag{18}
\]

In particular,

\[
I_{k-1,j}\subset I_{k,j}\subset K_{k,j}.
\tag{19}
\]

These inclusions also verify inductively the claimed lower bound on gaps
between different clusters. Because \(I_{k,j}\subset K_{k,j}\), (14)
proves the mesh bound (10); the confinement (15) proves its other
inclusion. The induction is complete.

The tolerance choices are not circular. The current spectral mesh is
fixed by increasing block lengths, without increasing the allowed
perturbation amplitude. The next step then uses \(\delta_k=a_{k+1}/8\)
to ensure (13).

\hypertarget{1031-the-limiting-spectrum-is-exactly-a-finite-interval-union}{%
\section{4. The limiting spectrum is exactly a finite interval
union}\label{1031-the-limiting-spectrum-is-exactly-a-finite-interval-union}}

The summable perturbation bounds imply that \(V_k\) converges uniformly
to a bounded real limit-periodic potential \(W\). Moreover,

\[
\|W-V\|_\infty
\leq\|W-V_0\|_\infty+\|V_0-V\|_\infty
<\frac t3+\frac{\varepsilon}{2}
<\varepsilon,
\tag{20}
\]

and

\[
\|W-V_k\|_\infty\leq\frac{a_k}{3}.
\tag{21}
\]

By (17)--(19), each sequence of left endpoints is decreasing and
bounded, and each sequence of right endpoints is increasing and bounded.
Write their limits as \(\ell_j,r_j\). The limiting intervals are
nondegenerate because they contain \(I_{0,j}\). Distinct intervals
remain separated by at least \(G-2t/3>0\). Thus

\[
K_\infty=\bigcup_{j=1}^{M}[\ell_j,r_j]
\tag{22}
\]

is a finite union of exactly \(M\) disjoint nondegenerate closed
intervals. The endpoint movement after stage \(k\) is at most
\(\sum_{s>k}a_s=a_k/3\), so

\[
d_H\!\left(\bigcup_j I_{k,j},K_\infty\right)
\leq\frac{a_k}{3}.
\tag{23}
\]

The mesh and confinement in (10) say

\[
d_H\!\left(S_k,\bigcup_jI_{k,j}\right)\leq\delta_k.
\tag{24}
\]

At the same time, norm continuity of the spectrum and (21) give

\[
d_H(S_k,\sigma(H_W))\leq\frac{a_k}{3}.
\tag{25}
\]

Combining (23)--(25) and letting \(k\to\infty\) yields

\[
\boxed{\sigma(H_W)=K_\infty.}
\tag{26}
\]

This proves the asserted conclusion for the entire spectrum. In
particular, if the initial periodic potential \(V_0\) is constant, its
spectrum is one interval and the construction retains \(M=1\).

\hypertarget{1031-the-limit-is-nonperiodic}{%
\section{5. The limit is nonperiodic}\label{1031-the-limit-is-nonperiodic}}

Fix an integer \(p\geq1\), and choose \(k\geq p\). The marker property
(11) gives a site \(n_{k,p}\) where the \(p\)-shift difference of
\(V_k\) has magnitude at least \(2a_k\). The future tail cannot erase
that difference:

\[
\begin{aligned}
|W(n_{k,p}+p)-W(n_{k,p})|
&\geq2a_k-2\|W-V_k\|_\infty\\
&\geq2a_k-\frac{2a_k}{3}
=\frac{4a_k}{3}>0.
\end{aligned}
\]

Therefore \(p\) is not a period of \(W\). Since \(p\) was arbitrary,
\(W\) is nonperiodic. Equations (20) and (26) now prove the theorem.

\hypertarget{1031-scope-and-references}{%
\section{Scope and references}\label{1031-scope-and-references}}

The construction guarantees that the entire spectrum is a finite union
of nondegenerate closed intervals. It can be a single interval when the
initial periodic approximant has connected spectrum; this includes
approximation of any constant potential. The argument makes no assertion
about the spectral type, so it does not resolve the separate mixed-type
or coupling-dependence questions.

\begin{enumerate}
\raggedright
\def\labelenumi{\arabic{enumi}.}
\tightlist
\item
  D. Damanik and J. Fillman, \emph{Spectral Properties of Limit-Periodic
  Operators}, arXiv:1802.05794v1 (2018), Question 10.4, printed p.~36;
  subsequently Chapter 15 in \emph{Analysis and Geometry on Graphs and
  Manifolds} (2020). The question explicitly requires density among
  nonperiodic limit-periodic potentials.
  \url{https://arxiv.org/pdf/1802.05794}.
\item
  D. Damanik, Z. Gan, and H. Krüger, \emph{Limit-periodic Schrödinger
  operators with a discontinuous Lyapunov exponent}, Journal of
  Functional Analysis \textbf{279} (2020), 108565. Its Theorem 1.3
  concerns density of discontinuous Lyapunov exponents and is not itself
  the spectral-interior statement proved here.
  \url{https://doi.org/10.1016/j.jfa.2020.108565}.
\end{enumerate}

The proof uses only elementary periodic approximation, the
approximate-eigenvector criterion, and norm continuity of the spectrum
for bounded self-adjoint operators. No bibliographic-priority claim is
made.
\end{document}
